\documentclass[11pt,a4paper]{article}

\usepackage[T1]{fontenc}
\usepackage[utf8]{inputenc}
\usepackage{lmodern}
\usepackage[margin=2.4cm]{geometry}
\usepackage{amsmath,amssymb}
\usepackage{booktabs}
\usepackage{longtable}
\usepackage{array}
\usepackage{graphicx}
\usepackage{pgfplots}
\pgfplotsset{compat=1.18}
\usepackage{microtype}
\usepackage[htt]{hyphenat}
\usepackage[hidelinks]{hyperref}
\usepackage{url}
\usepackage[numbers,sort&compress]{natbib}
\setlength{\emergencystretch}{2em}

% Generated by `papertool ledger` from data/measurements.json. A \val lookup with no macro
% prints a visible marker; `papertool numbers` fails the build on one.
% Generated by `papertool ledger` from data/measurements.json -- do not edit.
% Every quantitative claim in the paper resolves through one of these names.
\makeatletter
\expandafter\def\csname pv@e01_contextuality.x\endcsname{1}
\expandafter\def\csname pv@e01_contextuality.peak_ws_mb\endcsname{27.02}
\expandafter\def\csname pv@e01_contextuality.peak_ws_sd_mb\endcsname{0}
\expandafter\def\csname pv@e01_contextuality.elapsed_s\endcsname{0.4169}
\expandafter\def\csname pv@e01_contextuality.mermin_contextual_fraction\endcsname{0.333}
\expandafter\def\csname pvci@e01_contextuality.mermin_contextual_fraction\endcsname{0.333 (0)}
\expandafter\def\csname pvsd@e01_contextuality.mermin_contextual_fraction\endcsname{0}
\expandafter\def\csname pvn@e01_contextuality.mermin_contextual_fraction\endcsname{1}
\expandafter\def\csname pvmin@e01_contextuality.mermin_contextual_fraction\endcsname{0.333}
\expandafter\def\csname pvmax@e01_contextuality.mermin_contextual_fraction\endcsname{0.333}
\expandafter\def\csname pv@e01_contextuality.mermin_holonomy\endcsname{1}
\expandafter\def\csname pvci@e01_contextuality.mermin_holonomy\endcsname{1 (0)}
\expandafter\def\csname pvsd@e01_contextuality.mermin_holonomy\endcsname{0}
\expandafter\def\csname pvn@e01_contextuality.mermin_holonomy\endcsname{1}
\expandafter\def\csname pvmin@e01_contextuality.mermin_holonomy\endcsname{1}
\expandafter\def\csname pvmax@e01_contextuality.mermin_holonomy\endcsname{1}
\expandafter\def\csname pv@e01_contextuality.mermin_maxsat\endcsname{2}
\expandafter\def\csname pvci@e01_contextuality.mermin_maxsat\endcsname{2 (0)}
\expandafter\def\csname pvsd@e01_contextuality.mermin_maxsat\endcsname{0}
\expandafter\def\csname pvn@e01_contextuality.mermin_maxsat\endcsname{1}
\expandafter\def\csname pvmin@e01_contextuality.mermin_maxsat\endcsname{2}
\expandafter\def\csname pvmax@e01_contextuality.mermin_maxsat\endcsname{2}
\expandafter\def\csname pv@e01_contextuality.stdout_sha256\endcsname{99c6e74aa724c6986fa99f9a4e91476a387e9ba2605ab5709b989945b810f741}
\expandafter\def\csname pvuniq@e01_contextuality.stdout_sha256\endcsname{1}
\expandafter\def\csname pv@e01_contextuality_repeat.x\endcsname{1}
\expandafter\def\csname pv@e01_contextuality_repeat.peak_ws_mb\endcsname{27.2}
\expandafter\def\csname pv@e01_contextuality_repeat.peak_ws_sd_mb\endcsname{0.0635}
\expandafter\def\csname pv@e01_contextuality_repeat.elapsed_s\endcsname{0.3043}
\expandafter\def\csname pv@e01_contextuality_repeat.stdout_sha256\endcsname{99c6e74aa724c6986fa99f9a4e91476a387e9ba2605ab5709b989945b810f741}
\expandafter\def\csname pvuniq@e01_contextuality_repeat.stdout_sha256\endcsname{1}
\expandafter\def\csname pv@e05_diagonal_lemma.x\endcsname{5}
\expandafter\def\csname pv@e05_diagonal_lemma.peak_ws_mb\endcsname{27.19}
\expandafter\def\csname pv@e05_diagonal_lemma.peak_ws_sd_mb\endcsname{0}
\expandafter\def\csname pv@e05_diagonal_lemma.elapsed_s\endcsname{1.347}
\expandafter\def\csname pv@e05_diagonal_lemma.selfgauge_total\endcsname{120}
\expandafter\def\csname pvci@e05_diagonal_lemma.selfgauge_total\endcsname{120 (0)}
\expandafter\def\csname pvsd@e05_diagonal_lemma.selfgauge_total\endcsname{0}
\expandafter\def\csname pvn@e05_diagonal_lemma.selfgauge_total\endcsname{1}
\expandafter\def\csname pvmin@e05_diagonal_lemma.selfgauge_total\endcsname{120}
\expandafter\def\csname pvmax@e05_diagonal_lemma.selfgauge_total\endcsname{120}
\expandafter\def\csname pv@e05_diagonal_lemma.selfgauge_trivial_fixed\endcsname{60}
\expandafter\def\csname pvci@e05_diagonal_lemma.selfgauge_trivial_fixed\endcsname{60 (0)}
\expandafter\def\csname pvsd@e05_diagonal_lemma.selfgauge_trivial_fixed\endcsname{0}
\expandafter\def\csname pvn@e05_diagonal_lemma.selfgauge_trivial_fixed\endcsname{1}
\expandafter\def\csname pvmin@e05_diagonal_lemma.selfgauge_trivial_fixed\endcsname{60}
\expandafter\def\csname pvmax@e05_diagonal_lemma.selfgauge_trivial_fixed\endcsname{60}
\expandafter\def\csname pv@e05_diagonal_lemma.stdout_sha256\endcsname{03573f7eb327a3d7961c525a3732f1729de3846671fdddacdfe810bf1c2b8f36}
\expandafter\def\csname pvuniq@e05_diagonal_lemma.stdout_sha256\endcsname{1}
\expandafter\def\csname pv@e06_frustration.x\endcsname{6}
\expandafter\def\csname pv@e06_frustration.peak_ws_mb\endcsname{34.32}
\expandafter\def\csname pv@e06_frustration.peak_ws_sd_mb\endcsname{0}
\expandafter\def\csname pv@e06_frustration.elapsed_s\endcsname{10.33}
\expandafter\def\csname pv@e06_frustration.stdout_sha256\endcsname{e3d5fda89af87d6f2e607e917f7d0389fe86362a1f4dffa1548fde102de173a9}
\expandafter\def\csname pvuniq@e06_frustration.stdout_sha256\endcsname{1}
\expandafter\def\csname pv@e11_repair_hardness.x\endcsname{11}
\expandafter\def\csname pv@e11_repair_hardness.peak_ws_mb\endcsname{31.42}
\expandafter\def\csname pv@e11_repair_hardness.peak_ws_sd_mb\endcsname{0}
\expandafter\def\csname pv@e11_repair_hardness.elapsed_s\endcsname{0.4128}
\expandafter\def\csname pv@e11_repair_hardness.frustration_v8_index\endcsname{3}
\expandafter\def\csname pvci@e11_repair_hardness.frustration_v8_index\endcsname{3 (0)}
\expandafter\def\csname pvsd@e11_repair_hardness.frustration_v8_index\endcsname{0}
\expandafter\def\csname pvn@e11_repair_hardness.frustration_v8_index\endcsname{1}
\expandafter\def\csname pvmin@e11_repair_hardness.frustration_v8_index\endcsname{3}
\expandafter\def\csname pvmax@e11_repair_hardness.frustration_v8_index\endcsname{3}
\expandafter\def\csname pv@e11_repair_hardness.frustration_v8_maxsat\endcsname{13}
\expandafter\def\csname pvci@e11_repair_hardness.frustration_v8_maxsat\endcsname{13 (0)}
\expandafter\def\csname pvsd@e11_repair_hardness.frustration_v8_maxsat\endcsname{0}
\expandafter\def\csname pvn@e11_repair_hardness.frustration_v8_maxsat\endcsname{1}
\expandafter\def\csname pvmin@e11_repair_hardness.frustration_v8_maxsat\endcsname{13}
\expandafter\def\csname pvmax@e11_repair_hardness.frustration_v8_maxsat\endcsname{13}
\expandafter\def\csname pv@e11_repair_hardness.stdout_sha256\endcsname{07c52c6f0a586897bd0efe25f483140e15cb147b25da0601f749dda0687b8e00}
\expandafter\def\csname pvuniq@e11_repair_hardness.stdout_sha256\endcsname{1}
\expandafter\def\csname pv@e03_topological_memory.x\endcsname{3}
\expandafter\def\csname pv@e03_topological_memory.peak_ws_mb\endcsname{31.45}
\expandafter\def\csname pv@e03_topological_memory.peak_ws_sd_mb\endcsname{0.1407}
\expandafter\def\csname pv@e03_topological_memory.elapsed_s\endcsname{0.4149}
\expandafter\def\csname pv@e03_topological_memory.noise05_hopfield\endcsname{0.377}
\expandafter\def\csname pvci@e03_topological_memory.noise05_hopfield\endcsname{0.377 (0)}
\expandafter\def\csname pvsd@e03_topological_memory.noise05_hopfield\endcsname{0}
\expandafter\def\csname pvn@e03_topological_memory.noise05_hopfield\endcsname{5}
\expandafter\def\csname pvmin@e03_topological_memory.noise05_hopfield\endcsname{0.377}
\expandafter\def\csname pvmax@e03_topological_memory.noise05_hopfield\endcsname{0.377}
\expandafter\def\csname pv@e03_topological_memory.noise05_parity_struct\endcsname{0.138}
\expandafter\def\csname pvci@e03_topological_memory.noise05_parity_struct\endcsname{0.138 (0)}
\expandafter\def\csname pvsd@e03_topological_memory.noise05_parity_struct\endcsname{0}
\expandafter\def\csname pvn@e03_topological_memory.noise05_parity_struct\endcsname{5}
\expandafter\def\csname pvmin@e03_topological_memory.noise05_parity_struct\endcsname{0.138}
\expandafter\def\csname pvmax@e03_topological_memory.noise05_parity_struct\endcsname{0.138}
\expandafter\def\csname pv@e03_topological_memory.noise05_parity_value\endcsname{0}
\expandafter\def\csname pvci@e03_topological_memory.noise05_parity_value\endcsname{0 (0)}
\expandafter\def\csname pvsd@e03_topological_memory.noise05_parity_value\endcsname{0}
\expandafter\def\csname pvn@e03_topological_memory.noise05_parity_value\endcsname{5}
\expandafter\def\csname pvmin@e03_topological_memory.noise05_parity_value\endcsname{0}
\expandafter\def\csname pvmax@e03_topological_memory.noise05_parity_value\endcsname{0}
\expandafter\def\csname pv@e03_topological_memory.noise10_hopfield\endcsname{0.897}
\expandafter\def\csname pvci@e03_topological_memory.noise10_hopfield\endcsname{0.897 (0)}
\expandafter\def\csname pvsd@e03_topological_memory.noise10_hopfield\endcsname{0}
\expandafter\def\csname pvn@e03_topological_memory.noise10_hopfield\endcsname{5}
\expandafter\def\csname pvmin@e03_topological_memory.noise10_hopfield\endcsname{0.897}
\expandafter\def\csname pvmax@e03_topological_memory.noise10_hopfield\endcsname{0.897}
\expandafter\def\csname pv@e03_topological_memory.noise10_parity_struct\endcsname{0.246}
\expandafter\def\csname pvci@e03_topological_memory.noise10_parity_struct\endcsname{0.246 (0)}
\expandafter\def\csname pvsd@e03_topological_memory.noise10_parity_struct\endcsname{0}
\expandafter\def\csname pvn@e03_topological_memory.noise10_parity_struct\endcsname{5}
\expandafter\def\csname pvmin@e03_topological_memory.noise10_parity_struct\endcsname{0.246}
\expandafter\def\csname pvmax@e03_topological_memory.noise10_parity_struct\endcsname{0.246}
\expandafter\def\csname pv@e03_topological_memory.noise10_parity_value\endcsname{0}
\expandafter\def\csname pvci@e03_topological_memory.noise10_parity_value\endcsname{0 (0)}
\expandafter\def\csname pvsd@e03_topological_memory.noise10_parity_value\endcsname{0}
\expandafter\def\csname pvn@e03_topological_memory.noise10_parity_value\endcsname{5}
\expandafter\def\csname pvmin@e03_topological_memory.noise10_parity_value\endcsname{0}
\expandafter\def\csname pvmax@e03_topological_memory.noise10_parity_value\endcsname{0}
\expandafter\def\csname pv@e03_topological_memory.noise20_hopfield\endcsname{1}
\expandafter\def\csname pvci@e03_topological_memory.noise20_hopfield\endcsname{1 (0)}
\expandafter\def\csname pvsd@e03_topological_memory.noise20_hopfield\endcsname{0}
\expandafter\def\csname pvn@e03_topological_memory.noise20_hopfield\endcsname{5}
\expandafter\def\csname pvmin@e03_topological_memory.noise20_hopfield\endcsname{1}
\expandafter\def\csname pvmax@e03_topological_memory.noise20_hopfield\endcsname{1}
\expandafter\def\csname pv@e03_topological_memory.noise20_parity_struct\endcsname{0.385}
\expandafter\def\csname pvci@e03_topological_memory.noise20_parity_struct\endcsname{0.385 (0)}
\expandafter\def\csname pvsd@e03_topological_memory.noise20_parity_struct\endcsname{0}
\expandafter\def\csname pvn@e03_topological_memory.noise20_parity_struct\endcsname{5}
\expandafter\def\csname pvmin@e03_topological_memory.noise20_parity_struct\endcsname{0.385}
\expandafter\def\csname pvmax@e03_topological_memory.noise20_parity_struct\endcsname{0.385}
\expandafter\def\csname pv@e03_topological_memory.noise30_hopfield\endcsname{1}
\expandafter\def\csname pvci@e03_topological_memory.noise30_hopfield\endcsname{1 (0)}
\expandafter\def\csname pvsd@e03_topological_memory.noise30_hopfield\endcsname{0}
\expandafter\def\csname pvn@e03_topological_memory.noise30_hopfield\endcsname{5}
\expandafter\def\csname pvmin@e03_topological_memory.noise30_hopfield\endcsname{1}
\expandafter\def\csname pvmax@e03_topological_memory.noise30_hopfield\endcsname{1}
\expandafter\def\csname pv@e03_topological_memory.noise30_parity_struct\endcsname{0.463}
\expandafter\def\csname pvci@e03_topological_memory.noise30_parity_struct\endcsname{0.463 (0)}
\expandafter\def\csname pvsd@e03_topological_memory.noise30_parity_struct\endcsname{0}
\expandafter\def\csname pvn@e03_topological_memory.noise30_parity_struct\endcsname{5}
\expandafter\def\csname pvmin@e03_topological_memory.noise30_parity_struct\endcsname{0.463}
\expandafter\def\csname pvmax@e03_topological_memory.noise30_parity_struct\endcsname{0.463}
\expandafter\def\csname pv@e03_topological_memory.noise40_hopfield\endcsname{1}
\expandafter\def\csname pvci@e03_topological_memory.noise40_hopfield\endcsname{1 (0)}
\expandafter\def\csname pvsd@e03_topological_memory.noise40_hopfield\endcsname{0}
\expandafter\def\csname pvn@e03_topological_memory.noise40_hopfield\endcsname{5}
\expandafter\def\csname pvmin@e03_topological_memory.noise40_hopfield\endcsname{1}
\expandafter\def\csname pvmax@e03_topological_memory.noise40_hopfield\endcsname{1}
\expandafter\def\csname pv@e03_topological_memory.noise40_parity_struct\endcsname{0.497}
\expandafter\def\csname pvci@e03_topological_memory.noise40_parity_struct\endcsname{0.497 (0)}
\expandafter\def\csname pvsd@e03_topological_memory.noise40_parity_struct\endcsname{0}
\expandafter\def\csname pvn@e03_topological_memory.noise40_parity_struct\endcsname{5}
\expandafter\def\csname pvmin@e03_topological_memory.noise40_parity_struct\endcsname{0.497}
\expandafter\def\csname pvmax@e03_topological_memory.noise40_parity_struct\endcsname{0.497}
\expandafter\def\csname pv@e03_topological_memory.stdout_sha256\endcsname{9bef752b39ae15d0fa31fbb094fbdc79bef00853a79724c4606518f5de6fce05}
\expandafter\def\csname pvuniq@e03_topological_memory.stdout_sha256\endcsname{1}
\expandafter\def\csname pv@e12_holonomy_code.x\endcsname{12}
\expandafter\def\csname pv@e12_holonomy_code.peak_ws_mb\endcsname{31.6}
\expandafter\def\csname pv@e12_holonomy_code.peak_ws_sd_mb\endcsname{0}
\expandafter\def\csname pv@e12_holonomy_code.elapsed_s\endcsname{1.102}
\expandafter\def\csname pv@e12_holonomy_code.code_detect_v7_p02\endcsname{87}
\expandafter\def\csname pvci@e12_holonomy_code.code_detect_v7_p02\endcsname{87 (0)}
\expandafter\def\csname pvsd@e12_holonomy_code.code_detect_v7_p02\endcsname{0}
\expandafter\def\csname pvn@e12_holonomy_code.code_detect_v7_p02\endcsname{1}
\expandafter\def\csname pvmin@e12_holonomy_code.code_detect_v7_p02\endcsname{87}
\expandafter\def\csname pvmax@e12_holonomy_code.code_detect_v7_p02\endcsname{87}
\expandafter\def\csname pv@e12_holonomy_code.code_localise_v7_p02\endcsname{64}
\expandafter\def\csname pvci@e12_holonomy_code.code_localise_v7_p02\endcsname{64 (0)}
\expandafter\def\csname pvsd@e12_holonomy_code.code_localise_v7_p02\endcsname{0}
\expandafter\def\csname pvn@e12_holonomy_code.code_localise_v7_p02\endcsname{1}
\expandafter\def\csname pvmin@e12_holonomy_code.code_localise_v7_p02\endcsname{64}
\expandafter\def\csname pvmax@e12_holonomy_code.code_localise_v7_p02\endcsname{64}
\expandafter\def\csname pv@e12_holonomy_code.code_rank_v7\endcsname{4.3}
\expandafter\def\csname pvci@e12_holonomy_code.code_rank_v7\endcsname{4.3 (0)}
\expandafter\def\csname pvsd@e12_holonomy_code.code_rank_v7\endcsname{0}
\expandafter\def\csname pvn@e12_holonomy_code.code_rank_v7\endcsname{1}
\expandafter\def\csname pvmin@e12_holonomy_code.code_rank_v7\endcsname{4.3}
\expandafter\def\csname pvmax@e12_holonomy_code.code_rank_v7\endcsname{4.3}
\expandafter\def\csname pv@e12_holonomy_code.code_undetectable_v7_p02\endcsname{40}
\expandafter\def\csname pvci@e12_holonomy_code.code_undetectable_v7_p02\endcsname{40 (0)}
\expandafter\def\csname pvsd@e12_holonomy_code.code_undetectable_v7_p02\endcsname{0}
\expandafter\def\csname pvn@e12_holonomy_code.code_undetectable_v7_p02\endcsname{1}
\expandafter\def\csname pvmin@e12_holonomy_code.code_undetectable_v7_p02\endcsname{40}
\expandafter\def\csname pvmax@e12_holonomy_code.code_undetectable_v7_p02\endcsname{40}
\expandafter\def\csname pv@e12_holonomy_code.stdout_sha256\endcsname{357851df89327aa0460a0e41809519f31d98c32a91688c3eb7e5ca57777682dd}
\expandafter\def\csname pvuniq@e12_holonomy_code.stdout_sha256\endcsname{1}
\expandafter\def\csname pv@e04_novelty.x\endcsname{4}
\expandafter\def\csname pv@e04_novelty.peak_ws_mb\endcsname{31.5}
\expandafter\def\csname pv@e04_novelty.peak_ws_sd_mb\endcsname{0}
\expandafter\def\csname pv@e04_novelty.elapsed_s\endcsname{11.93}
\expandafter\def\csname pv@e04_novelty.novelty_mu002\endcsname{14}
\expandafter\def\csname pvci@e04_novelty.novelty_mu002\endcsname{14 (0)}
\expandafter\def\csname pvsd@e04_novelty.novelty_mu002\endcsname{0}
\expandafter\def\csname pvn@e04_novelty.novelty_mu002\endcsname{1}
\expandafter\def\csname pvmin@e04_novelty.novelty_mu002\endcsname{14}
\expandafter\def\csname pvmax@e04_novelty.novelty_mu002\endcsname{14}
\expandafter\def\csname pv@e04_novelty.novelty_mu01\endcsname{66}
\expandafter\def\csname pvci@e04_novelty.novelty_mu01\endcsname{66 (0)}
\expandafter\def\csname pvsd@e04_novelty.novelty_mu01\endcsname{0}
\expandafter\def\csname pvn@e04_novelty.novelty_mu01\endcsname{1}
\expandafter\def\csname pvmin@e04_novelty.novelty_mu01\endcsname{66}
\expandafter\def\csname pvmax@e04_novelty.novelty_mu01\endcsname{66}
\expandafter\def\csname pv@e04_novelty.novelty_mu05\endcsname{314}
\expandafter\def\csname pvci@e04_novelty.novelty_mu05\endcsname{314 (0)}
\expandafter\def\csname pvsd@e04_novelty.novelty_mu05\endcsname{0}
\expandafter\def\csname pvn@e04_novelty.novelty_mu05\endcsname{1}
\expandafter\def\csname pvmin@e04_novelty.novelty_mu05\endcsname{314}
\expandafter\def\csname pvmax@e04_novelty.novelty_mu05\endcsname{314}
\expandafter\def\csname pv@e04_novelty.novelty_value_noise\endcsname{0}
\expandafter\def\csname pvci@e04_novelty.novelty_value_noise\endcsname{0 (0)}
\expandafter\def\csname pvsd@e04_novelty.novelty_value_noise\endcsname{0}
\expandafter\def\csname pvn@e04_novelty.novelty_value_noise\endcsname{1}
\expandafter\def\csname pvmin@e04_novelty.novelty_value_noise\endcsname{0}
\expandafter\def\csname pvmax@e04_novelty.novelty_value_noise\endcsname{0}
\expandafter\def\csname pv@e04_novelty.stdout_sha256\endcsname{50947abcd41213372b1bfddf96bf10e15218ae84ef971ef9a9a2ad37e4348d65}
\expandafter\def\csname pvuniq@e04_novelty.stdout_sha256\endcsname{1}
\expandafter\def\csname pv@e02_defect_thermodynamics.x\endcsname{2}
\expandafter\def\csname pv@e02_defect_thermodynamics.peak_ws_mb\endcsname{32.71}
\expandafter\def\csname pv@e02_defect_thermodynamics.peak_ws_sd_mb\endcsname{0}
\expandafter\def\csname pv@e02_defect_thermodynamics.elapsed_s\endcsname{250.9}
\expandafter\def\csname pv@e02_defect_thermodynamics.defect_events_mu0\endcsname{0}
\expandafter\def\csname pvci@e02_defect_thermodynamics.defect_events_mu0\endcsname{0 (0)}
\expandafter\def\csname pvsd@e02_defect_thermodynamics.defect_events_mu0\endcsname{0}
\expandafter\def\csname pvn@e02_defect_thermodynamics.defect_events_mu0\endcsname{1}
\expandafter\def\csname pvmin@e02_defect_thermodynamics.defect_events_mu0\endcsname{0}
\expandafter\def\csname pvmax@e02_defect_thermodynamics.defect_events_mu0\endcsname{0}
\expandafter\def\csname pv@e02_defect_thermodynamics.defect_events_mu002\endcsname{0.69}
\expandafter\def\csname pvci@e02_defect_thermodynamics.defect_events_mu002\endcsname{0.69 (0)}
\expandafter\def\csname pvsd@e02_defect_thermodynamics.defect_events_mu002\endcsname{0}
\expandafter\def\csname pvn@e02_defect_thermodynamics.defect_events_mu002\endcsname{1}
\expandafter\def\csname pvmin@e02_defect_thermodynamics.defect_events_mu002\endcsname{0.69}
\expandafter\def\csname pvmax@e02_defect_thermodynamics.defect_events_mu002\endcsname{0.69}
\expandafter\def\csname pv@e02_defect_thermodynamics.defect_events_mu01\endcsname{0.71}
\expandafter\def\csname pvci@e02_defect_thermodynamics.defect_events_mu01\endcsname{0.71 (0)}
\expandafter\def\csname pvsd@e02_defect_thermodynamics.defect_events_mu01\endcsname{0}
\expandafter\def\csname pvn@e02_defect_thermodynamics.defect_events_mu01\endcsname{1}
\expandafter\def\csname pvmin@e02_defect_thermodynamics.defect_events_mu01\endcsname{0.71}
\expandafter\def\csname pvmax@e02_defect_thermodynamics.defect_events_mu01\endcsname{0.71}
\expandafter\def\csname pv@e02_defect_thermodynamics.defect_events_mu02\endcsname{0.72}
\expandafter\def\csname pvci@e02_defect_thermodynamics.defect_events_mu02\endcsname{0.72 (0)}
\expandafter\def\csname pvsd@e02_defect_thermodynamics.defect_events_mu02\endcsname{0}
\expandafter\def\csname pvn@e02_defect_thermodynamics.defect_events_mu02\endcsname{1}
\expandafter\def\csname pvmin@e02_defect_thermodynamics.defect_events_mu02\endcsname{0.72}
\expandafter\def\csname pvmax@e02_defect_thermodynamics.defect_events_mu02\endcsname{0.72}
\expandafter\def\csname pv@e02_defect_thermodynamics.defect_events_mu05\endcsname{0.73}
\expandafter\def\csname pvci@e02_defect_thermodynamics.defect_events_mu05\endcsname{0.73 (0)}
\expandafter\def\csname pvsd@e02_defect_thermodynamics.defect_events_mu05\endcsname{0}
\expandafter\def\csname pvn@e02_defect_thermodynamics.defect_events_mu05\endcsname{1}
\expandafter\def\csname pvmin@e02_defect_thermodynamics.defect_events_mu05\endcsname{0.73}
\expandafter\def\csname pvmax@e02_defect_thermodynamics.defect_events_mu05\endcsname{0.73}
\expandafter\def\csname pv@e02_defect_thermodynamics.defect_hurst_mu0\endcsname{0.092}
\expandafter\def\csname pvci@e02_defect_thermodynamics.defect_hurst_mu0\endcsname{0.092 (0)}
\expandafter\def\csname pvsd@e02_defect_thermodynamics.defect_hurst_mu0\endcsname{0}
\expandafter\def\csname pvn@e02_defect_thermodynamics.defect_hurst_mu0\endcsname{1}
\expandafter\def\csname pvmin@e02_defect_thermodynamics.defect_hurst_mu0\endcsname{0.092}
\expandafter\def\csname pvmax@e02_defect_thermodynamics.defect_hurst_mu0\endcsname{0.092}
\expandafter\def\csname pv@e02_defect_thermodynamics.defect_rho_mu0\endcsname{0.7}
\expandafter\def\csname pvci@e02_defect_thermodynamics.defect_rho_mu0\endcsname{0.7 (0)}
\expandafter\def\csname pvsd@e02_defect_thermodynamics.defect_rho_mu0\endcsname{0}
\expandafter\def\csname pvn@e02_defect_thermodynamics.defect_rho_mu0\endcsname{1}
\expandafter\def\csname pvmin@e02_defect_thermodynamics.defect_rho_mu0\endcsname{0.7}
\expandafter\def\csname pvmax@e02_defect_thermodynamics.defect_rho_mu0\endcsname{0.7}
\expandafter\def\csname pv@e02_defect_thermodynamics.defect_rho_mu02\endcsname{0.72}
\expandafter\def\csname pvci@e02_defect_thermodynamics.defect_rho_mu02\endcsname{0.72 (0)}
\expandafter\def\csname pvsd@e02_defect_thermodynamics.defect_rho_mu02\endcsname{0}
\expandafter\def\csname pvn@e02_defect_thermodynamics.defect_rho_mu02\endcsname{1}
\expandafter\def\csname pvmin@e02_defect_thermodynamics.defect_rho_mu02\endcsname{0.72}
\expandafter\def\csname pvmax@e02_defect_thermodynamics.defect_rho_mu02\endcsname{0.72}
\expandafter\def\csname pv@e02_defect_thermodynamics.stdout_sha256\endcsname{333752a6f872d3f439642f2dc05765f3cc2426b5c6db3fddda7edd08ab1c6f87}
\expandafter\def\csname pvuniq@e02_defect_thermodynamics.stdout_sha256\endcsname{1}
\expandafter\def\csname pv@selfgauge_fraction.value\endcsname{0.5}
\expandafter\def\csname pvci@selfgauge_fraction.value\endcsname{0.5 (0)}
\expandafter\def\csname pvsd@selfgauge_fraction.value\endcsname{0}
\expandafter\def\csname pvn@selfgauge_fraction.value\endcsname{1}
\expandafter\def\csname pvmin@selfgauge_fraction.value\endcsname{0.5}
\expandafter\def\csname pvmax@selfgauge_fraction.value\endcsname{0.5}
\expandafter\def\csname pv@noise10_hopfield_over_parity.value\endcsname{3.646}
\expandafter\def\csname pvci@noise10_hopfield_over_parity.value\endcsname{3.646 (0)}
\expandafter\def\csname pvsd@noise10_hopfield_over_parity.value\endcsname{0}
\expandafter\def\csname pvn@noise10_hopfield_over_parity.value\endcsname{1}
\expandafter\def\csname pvmin@noise10_hopfield_over_parity.value\endcsname{3.646}
\expandafter\def\csname pvmax@noise10_hopfield_over_parity.value\endcsname{3.646}
\expandafter\def\csname pv@mcerr_parity_struct.value\endcsname{0.0068}
\expandafter\def\csname pvci@mcerr_parity_struct.value\endcsname{0.0068 (0)}
\expandafter\def\csname pvsd@mcerr_parity_struct.value\endcsname{0}
\expandafter\def\csname pvn@mcerr_parity_struct.value\endcsname{1}
\expandafter\def\csname pvmin@mcerr_parity_struct.value\endcsname{0.0068}
\expandafter\def\csname pvmax@mcerr_parity_struct.value\endcsname{0.0068}
\expandafter\def\csname pv@novelty_structure_over_value.value\endcsname{14}
\expandafter\def\csname pvci@novelty_structure_over_value.value\endcsname{14 (0)}
\expandafter\def\csname pvsd@novelty_structure_over_value.value\endcsname{0}
\expandafter\def\csname pvn@novelty_structure_over_value.value\endcsname{1}
\expandafter\def\csname pvmin@novelty_structure_over_value.value\endcsname{14}
\expandafter\def\csname pvmax@novelty_structure_over_value.value\endcsname{14}
\makeatother

\newcommand{\val}[1]{\ifcsname pv@#1\endcsname\csname pv@#1\endcsname\else\textbf{??#1??}\fi}
\newcommand{\valci}[1]{\ifcsname pvci@#1\endcsname\csname pvci@#1\endcsname\else\textbf{??#1??}\fi}
\newcommand{\valsd}[1]{\ifcsname pvsd@#1\endcsname\csname pvsd@#1\endcsname\else\textbf{??#1??}\fi}
\newcommand{\valn}[1]{\ifcsname pvn@#1\endcsname\csname pvn@#1\endcsname\else\textbf{??#1??}\fi}
\newcommand{\valfrac}[1]{\ifcsname pvfrac@#1\endcsname\csname pvfrac@#1\endcsname\else\textbf{??#1??}\fi}
\newcommand{\valtrue}[1]{\ifcsname pvtrue@#1\endcsname\csname pvtrue@#1\endcsname\else\textbf{??#1??}\fi}

\title{\bfseries Parity holonomy in Boolean constraint networks\\[2pt]
\large A reproduction study of five claims, one refutation, and a checksum that is not a memory}
\author{Konstantin Parashchuk\\\small Independent researcher, Saint Petersburg, Russia\\
\small \texttt{github.com/\allowbreak{}kostyk348}}
\date{17 November 2026}

\begin{document}
\maketitle

\begin{abstract}
\noindent
A repository sets out to test whether computation can be founded on paradox rather than on
state and rule, formalises the question over GF(2), turns every assertion into an experiment, and
reports its own results including the ones that failed. This paper reproduces that study
independently. All eight of its experiments are re-run from a pinned commit and compared against the
result files committed with the code: every one reproduces byte for byte, including the one that
draws four thousand random trials, which turns out to be deterministically seeded. Each of the five claims
is checked against numbers. The identification of paradox with contextuality holds: on a
negation-labelled cycle the invariant is non-trivial exactly when no global assignment exists,
and the Mermin triangle returns the expected obstruction. The refutation of the author's own
diagonal lemma holds as well: \val{e05_diagonal_lemma.selfgauge_trivial_fixed} of \val{e05_diagonal_lemma.selfgauge_total} self-gauge
systems admit a consistent fixed point, so self-reference does not force paradox. The third
claim needs a sharper statement than the repository gives it. Parity survives value noise
completely and structural noise partially (\val{e03_topological_memory.noise10_parity_struct} of
the bit is lost at 0.10 structural error, against
\val{e03_topological_memory.noise10_hopfield} for the associative memory it is compared with), but
as an error-correcting object it has minimum distance one: at \val{e12_holonomy_code.code_undetectable_v7_p02}
single edge errors, each placed on a bridge, are invisible to it. That makes it a checksum, which is what the
author says in a caveat and what this paper states as a theorem-shaped fact with a number
attached. Two claims are not reproduced here, for a reason the paper gives explicitly.
\end{abstract}

\noindent\textbf{Keywords:} parity holonomy; Boolean constraint networks; contextuality; error
detecting codes; frustration; reproducibility

\section{Introduction}

Most research repositories publish a claim and a demonstration. A smaller number publish a
claim, a demonstration, and a verdict on their own claim, including the verdicts that go against
the author. The subject of this paper is in the second group: \texttt{paradox-logic}
\citep{paradoxlogic} asks whether a computational system can be built on paradox as a foundation
rather than treating paradox as an error, formalises the question over GF(2), and then reports,
claim by claim, which parts of the construction survived.

The author's short answer, printed at the top of the README, is that the only non-trivial
structure in the whole construction is parity: a paradox is a non-trivial class in the first
cohomology of the network's negation labelling with $\mathbb{Z}/2$ coefficients, which makes the
invariant a useful coordinate and not a new class of computation. That is an unusually deflationary
conclusion to publish about one's own programme, and it is the reason the repository is worth
auditing: a deflationary claim is one that can be checked.

This paper does three things. It reproduces the experiments from a pinned commit and reports
whether the committed result files reappear exactly (\S\ref{sec:repro}). It checks the numbers
behind four of the five verdicts against the claims those numbers are supposed to support
(\S\ref{sec:claims}). And it sharpens the one verdict that is stated too weakly: the parity
invariant is described as "a checksum, not a memory", which is true, and this paper makes the
sense in which it is true precise by measuring the invariant's minimum distance
(\S\ref{sec:code}).

\section{Subject and its own verdicts}
\label{sec:subject}

The subject is \texttt{kostyk348/\allowbreak{}paradox-logic} at commit \texttt{02c66ff}. It is a Python
package with two layers: an algebra layer that treats a process network as a labelled graph over
a group and answers consistency questions by cohomology, and a learning layer that embeds finite
semigroups into networks, the decomposition tradition of \citet{krohn1965} with an $L^*$
state-discovery front end \citep{angluin1987}. Alongside the library sit 90 numbered experiment scripts and a
\texttt{results/\allowbreak{}} directory holding the printed output of each, committed next to the code that
produced it.

The README presents five directions with verdicts, and because the verdicts are the thing being
audited they are quoted here rather than paraphrased
(Table~\ref{tab:verdicts}).

\begin{table}[t]
\centering
\footnotesize
\begin{tabular}{@{}p{0.40\textwidth}p{0.22\textwidth}p{0.30\textwidth}@{}}
\toprule
direction as stated in the repository & its verdict & this paper \\
\midrule
Paradox is Mermin/GHZ contextuality & confirmed & reproduced, \S\ref{sec:claims} \\
Defect thermodynamics: criticality from paradox birth and death & not found & reproduced: \S\ref{sec:defects} \\
Topological (parity) memory & confirmed, "it is a checksum, not a memory" & reproduced and sharpened, \S\ref{sec:code} \\
Novelty as a change of topological class & core yes, criticality signal no & reproduced, \S\ref{sec:novelty} \\
Cohomological diagonal lemma: self-reference implies paradox & refuted & reproduced, \S\ref{sec:claims} \\
\bottomrule
\end{tabular}
\caption{The five verdicts the audited repository states for itself, and what this paper did with
each. Every verdict is now checked; the two the repository itself marks as negative or partial are
reproduced as such.}
\label{tab:verdicts}
\end{table}

The formal object is simple enough to state in one paragraph, which is a virtue for an audit. A
network is a graph whose edges carry a negation bit; a global assignment of values to vertices is
consistent when every edge's parity constraint is satisfied. Only the parity of the edge labels
around each cycle matters, so the whole question lives in the cycle space, and the obstruction is
a class in $H^1(G;\mathbb{Z}/2)$ of the labelling. The cohomological reading of contextuality is
standard \citep{abramsky2011,abramsky2012}; what the audited repository adds is a set of finite
experiments that make the correspondence checkable in a few seconds.

\section{Method}
\label{sec:method}

\paragraph{Plan first, then compare, then check numbers.} Every run is declared in
\texttt{experiments/\allowbreak{}plan.toml}: the script, its working directory, how many repeats, and what the
run is for. The harness records the child's stdout, hashes it, and compares the whole file
against the committed result with \texttt{cmp} semantics, so "reproduced" means byte for byte
rather than "prints similar numbers".

\paragraph{Where the numbers in this paper come from.} \texttt{papertool measure} turns the raw
output into metrics (explicit extractions for the printed tables) and aggregates them into
\texttt{data/\allowbreak{}measurements.json}; \texttt{papertool ledger} turns that into the \TeX{} macros used
in the text and a ledger row per value; \texttt{papertool numbers} fails the build if a lookup
does not resolve or if a decimal appears in the prose that is not a measurement or an explained
constant. The rule is the one \citet{sandve2013} argue for, enforced by a program instead of by
discipline.

\paragraph{The one statistical device.} A printed rate from a finite Monte Carlo run has a
sampling error whether or not the author prints it. For the structural-noise series the audited
script runs 4000 trials per
row on a single RNG stream; the binomial standard error at the printed
\val{e03_topological_memory.noise10_parity_struct} rate is $\pm\val{mcerr_parity_struct.value}$, which is the
resolution of the three printed digits and is quoted with it in \S\ref{sec:code}.

\section{Reproduction}
\label{sec:repro}

Eight experiments were re-run from the pinned commit: the contextuality enumeration, the
topological-memory comparison, the novelty scan, the diagonal-lemma enumeration, the frustration
control, the repair-hardness measurement and the holonomy-code measurement. All eight produced output identical
to the \texttt{results/\allowbreak{}} file committed with the repository, and the two exhaustive experiments
were run three times and once more respectively with the same result each time
(Table~\ref{tab:repro}).

\begin{table}[t]
\centering
\small
\begin{tabular}{llrl}
\toprule
experiment & what it enumerates & runs & identical to \texttt{results/\allowbreak{}} \\
\midrule
\texttt{01\_\allowbreak{}contextuality} & cycles $L=3\ldots6$, all label counts & 3 & yes \\
\texttt{03\_\allowbreak{}topological\_\allowbreak{}memory} & 5 noise levels, 3 dynamics & \valn{e03_topological_memory.noise10_hopfield} & yes \\
\texttt{05\_\allowbreak{}diagonal\_\allowbreak{}lemma} & $2^L$ label vectors, $L=3\ldots6$ & \valn{e05_diagonal_lemma.selfgauge_total} & yes \\
\texttt{06\_\allowbreak{}frustration} & spin glass at 8 temperatures & 1 & yes \\
\texttt{11\_\allowbreak{}repair\_\allowbreak{}hardness} & frustrated graphs $V=8\ldots14$ & 1 & yes \\
\texttt{12\_\allowbreak{}holonomy\_\allowbreak{}code} & random $S_3$-labelled graphs & 1 & yes \\
\bottomrule
\end{tabular}
\caption{Reproduction of the committed results. ``Runs'' is the number of processes; every one of
them produced output byte-identical to the corresponding file in \texttt{results/\allowbreak{}}, which is why
``identical'' is a single column and not a percentage.}
\label{tab:repro}
\end{table}

One detail is worth recording because it changes how the numbers should be read. The
topological-memory experiment looks stochastic and is not: it opens
\texttt{numpy.random.default\_\allowbreak{}rng(3)}, so five repeats gave identical digits, and the standard
deviation across repeats is exactly zero. The same is true of the holonomy-code and frustration
experiments. The good consequence is that the repository is reproducible in the strong sense; the
consequence a reader should not miss is that the Monte Carlo error is fixed and unquantified by
the code, which is why \S\ref{sec:code} attaches a binomial standard error to the printed rates.

\section{The claims, checked}
\label{sec:claims}

%%NUMBERS-STRICT-BEGIN
\subsection{Paradox is contextuality}

The first claim is the one the whole construction rests on, and it holds in the simplest possible
form. On a cycle with an odd number of negation labels there is no global assignment: the product
of the constraints around the cycle is $-1$ while the product of the labels is $+1$, and the
holonomy is therefore non-trivial. The re-run enumerates cycles from three to six vertices with
every label count and reports the invariant as \val{e01_contextuality.mermin_holonomy} with a
maximum satisfiable set of \val{e01_contextuality.mermin_maxsat} clauses of three for the Mermin
triangle, exactly the parity obstruction of \citet{mermin1990,mermin1993}. Nothing about this is
new mathematics, which is the point: the audit confirms that the implementation computes the
object the text says it computes.

\subsection{The diagonal lemma does not hold in its simple form}

The fifth claim is a refutation of the author's own hypothesis, and this is where a repository of
this kind earns its keep. The hypothesis was that a network whose labels depend on its own
holonomy class must be paradoxical. The test is exhaustive: for every base vector over
$L=3\ldots6$ vertices, apply the self-gauge map and look for a fixed point with trivial class.
The fixed point exists for \val{e05_diagonal_lemma.selfgauge_trivial_fixed} of \val{e05_diagonal_lemma.selfgauge_total} vectors, a
fraction of \val{selfgauge_fraction.value}. Self-reference does not force paradox; it forces it
about as often as a fair coin would.

\subsection{A negative result, reproduced: defects do not go critical}
\label{sec:defects}

The second claim is the one the repository reports as a failure, and the re-run reproduces both the
numbers and the failure. A structure-noise scan at sixteen vertices, four thousand steps and thirty
seeds reports the stationary defect density at five noise rates: \val{e02_defect_thermodynamics.defect_rho_mu0} at zero noise,
then \val{e02_defect_thermodynamics.defect_rho_mu002}, \val{e02_defect_thermodynamics.defect_rho_mu01}, \val{e02_defect_thermodynamics.defect_rho_mu05} and \val{e02_defect_thermodynamics.defect_rho_mu02} as the rate rises. The density does not
move. What moves is the count of class-change events: \val{e02_defect_thermodynamics.defect_events_mu0} at zero noise, then
\val{e02_defect_thermodynamics.defect_events_mu002}, \val{e02_defect_thermodynamics.defect_events_mu01}, \val{e02_defect_thermodynamics.defect_events_mu05} and \val{e02_defect_thermodynamics.defect_events_mu02}. The counts grow with the noise rate, which
is what Poisson counting does, not what an avalanche does. The Hurst exponent at zero noise is
\val{e02_defect_thermodynamics.defect_hurst_mu0}, far from the 0.5 of an uncorrelated series and nowhere near the range above it that
a critical process would need.

The audited script prints its own verdict in one line: the density is stationary, events scale with
the noise rate, there is no avalanche and no critical point. This audit confirms that negative
result rather than overturning it. A construction built on defects does not, at this size, drive
the defects into a critical state, and the repository's summary of its own data is the correct
reading of that data. Publishing the failure is the part of the audited work that deserves the
most credit, and reproducing it is the cheapest way to say so.

\subsection{Topological memory: what the numbers actually show}
\label{sec:memory}

The third claim compares two ways of holding a bit under noise. In the parity construction the bit
is the holonomy of an odd cycle, so flipping the \emph{values} at the vertices cannot touch it,
while flipping the \emph{labels} can. The baseline is a Hopfield network, which stores the bit in
vertex values and therefore degrades as the values do \citep{hopfield1982}.

The re-run reproduces the qualitative picture exactly (Table~\ref{tab:noise},
Figure~\ref{fig:noise}). Parity error under value noise is \val{e03_topological_memory.noise10_parity_value}
at every noise level tried, because the invariant does not read the values at all. Under
structural noise the parity bit degrades roughly linearly, from
\val{e03_topological_memory.noise05_parity_struct} at
0.05 to \val{e03_topological_memory.noise40_parity_struct} at
0.40. The Hopfield baseline is at
\val{e03_topological_memory.noise10_hopfield} at the same noise level where parity is at
\val{e03_topological_memory.noise10_parity_struct}, a factor of
\val{noise10_hopfield_over_parity.value}, and it saturates at
\val{e03_topological_memory.noise20_hopfield} by 0.20.

\begin{table}[t]
\centering
\small
\begin{tabular}{lrrr}
\toprule
noise per landmark & parity, value noise & parity, structural noise & Hopfield, value noise \\
\midrule
0.05 & \val{e03_topological_memory.noise05_parity_value} & \val{e03_topological_memory.noise05_parity_struct} & \val{e03_topological_memory.noise05_hopfield} \\
0.10 & \val{e03_topological_memory.noise10_parity_value} & \val{e03_topological_memory.noise10_parity_struct} & \val{e03_topological_memory.noise10_hopfield} \\
0.20     & \val{e03_topological_memory.noise10_parity_value} & \val{e03_topological_memory.noise20_parity_struct} & \val{e03_topological_memory.noise20_hopfield} \\
0.30     & \val{e03_topological_memory.noise10_parity_value} & \val{e03_topological_memory.noise30_parity_struct} & \val{e03_topological_memory.noise30_hopfield} \\
0.40     & \val{e03_topological_memory.noise10_parity_value} & \val{e03_topological_memory.noise40_parity_struct} & \val{e03_topological_memory.noise40_hopfield} \\
\bottomrule
\end{tabular}
\caption{Error rate of a stored bit under three dynamics, as printed by the audited experiment and
reproduced here. The second column is zero at every level because the parity invariant does not
read vertex values; the third column is where the bit can actually be lost. The Hopfield column
saturates. Each rate comes from one seeded stream of four thousand trials; the binomial standard
error at the middle row is $\pm\val{mcerr_parity_struct.value}$.}
\label{tab:noise}
\end{table}

\begin{figure}[t]
\centering
\begin{tikzpicture}
\begin{axis}[
  width=0.74\textwidth, height=5.2cm,
  xlabel={structural / value noise per landmark},
  ylabel={error rate of the stored bit},
  xmin=0, xmax=0.45, ymin=-0.05, ymax=1.1,
  grid=major, grid style={gray!25},
  legend pos=north west, legend style={font=\footnotesize},
]
\addplot+[mark=square*, thick] coordinates {
  (0.05,\val{e03_topological_memory.noise05_parity_struct})
  (0.10,\val{e03_topological_memory.noise10_parity_struct})
  (0.20,\val{e03_topological_memory.noise20_parity_struct})
  (0.30,\val{e03_topological_memory.noise30_parity_struct})
  (0.40,\val{e03_topological_memory.noise40_parity_struct})};
\addlegendentry{parity, structural noise}
\addplot+[mark=*, thick] coordinates {
  (0.05,\val{e03_topological_memory.noise05_hopfield})
  (0.10,\val{e03_topological_memory.noise10_hopfield})
  (0.20,\val{e03_topological_memory.noise20_hopfield})
  (0.30,\val{e03_topological_memory.noise30_hopfield})
  (0.40,\val{e03_topological_memory.noise40_hopfield})};
\addlegendentry{Hopfield, value noise}
\addplot+[mark=triangle*, thick] coordinates {
  (0.05,0) (0.10,0) (0.20,0) (0.30,0) (0.40,0)};
\addlegendentry{parity, value noise}
\end{axis}
\end{tikzpicture}
\caption{The three dynamics of Table~\ref{tab:noise}. The flat line at zero is not a triumph of
the method over noise in general; it is what happens when the noise is applied to a quantity the
invariant does not read.}
\label{fig:noise}
\end{figure}

Read carefully, this table says that parity is a robust \emph{detector} of one kind of damage and
a fragile \emph{store} of another, which is what the repository's caveat says in words. It does
not say the invariant is a memory, and the next section measures how far from a memory it is.

\section{The invariant as a code: distance one}
\label{sec:code}

The repository contains a further experiment that treats a consistent labelling as a codeword and
a corrupted edge as an error, and reports detection, localisation, and one more column: single
errors placed on a bridge, which no cycle passes through and which therefore leave every holonomy
unchanged. The re-run reproduces its output exactly, and the number that matters is
\val{e12_holonomy_code.code_undetectable_v7_p02} such errors at seven vertices and edge
probability 0.2 (Table~\ref{tab:code}).

\begin{table}[t]
\centering
\small
\begin{tabular}{lrrr}
\toprule
graph & cycle rank & detected & single errors invisible \\
\midrule
7 vertices, $p=0.2$ & \val{e12_holonomy_code.code_rank_v7} & \val{e12_holonomy_code.code_detect_v7_p02}\% & \val{e12_holonomy_code.code_undetectable_v7_p02} \\
\bottomrule
\end{tabular}
\caption{The holonomy code at one configuration, reproduced from the audited experiment. An error
on a bridge lies in no cycle, so it changes no holonomy: the minimum weight of an undetectable
error is one. Every number in the row is a value the audited experiment prints and this
package extracts; the row is one configuration, not a sweep.}
\label{tab:code}
\end{table}

A code whose minimum distance is one detects nothing and corrects nothing. The relevant comparison
is with a real topological memory, where the invariant has a distance that grows with the system
size and a local error model cannot cross it \citep{kitaev2003,dennis2002}. The parity invariant has
no such distance: every edge that is not on a cycle is invisible to it, and the minimum-weight
undetectable error is a single edge. That is exactly the honest content of the phrase "a checksum,
not a memory" \citep{hamming1950}, and it is stronger than a caveat because it is a number.

The same section of the repository measures what repair would cost, and the answer is the reason
the construction cannot be made into a memory by fiat. Detecting an inconsistency is one GF(2)
rank computation; repairing it minimally is MAX-CUT, which is NP-hard \citep{karp1972}. The
reproduced numbers make the gap concrete at eight vertices: a maximum satisfiable set of
\val{e11_repair_hardness.frustration_v8_maxsat} constraints and a frustration index of
\val{e11_repair_hardness.frustration_v8_index}. The frustration language is the standard one for
signed networks \citep{vannimenus1977}, and the control experiment on a spin glass
\citep{binder1986,goles1980} shows the transition the defect-gas construction did not produce.
%%NUMBERS-STRICT-END

\subsection{Novelty: orthogonal to value noise, linear in structural damage}
\label{sec:novelty}

The fourth claim is about novelty: does a change of topological class give a usable definition of
a new thing happening, as opposed to mere randomness? The re-run reproduces the printed answer
exactly. Under value noise the number of class changes is \val{e04_novelty.novelty_value_noise}, which is to say the
construction cannot see value randomness at all. Under structural noise the count grows with the
damage rate: \val{e04_novelty.novelty_mu002} class changes at a structural rate of 0.002, \val{e04_novelty.novelty_mu01} at 0.01 and
\val{e04_novelty.novelty_mu05} at 0.05.

That is the claim the repository marks as half-supported, and the numbers show why the
reservation is the right one. The construction delivers the first half (novelty is not noise, and
the two kinds of perturbation separate cleanly) and not the second (a criticality signal, which
is the part that would make it a mechanism rather than a detector). The count is roughly linear
in the structural rate over the range tested, which is consistent with a detector and not with a
transition.

\section{What this paper does not establish}
\label{sec:not}

\paragraph{Nothing was left out, and one exclusion was a mistake.} The first pass of this audit
excluded two experiments as too slow. The novelty scan turned out to run in minutes and the defect
scan in about five, so both are now in the plan and in the results; the one that stayed excluded
longest is the one whose numbers this paper quotes most carefully, because "runtime budget" was a
guess rather than a measurement. The CHANGELOG records the correction.

\paragraph{A claim-verdict format borrowed from replication practice.} The way this paper reports
results, \citep{osc2015} one verdict per claim with the number behind it and a separate list of
what was not checked, follows the replication literature rather than the repository being audited
\citep{ioannidis2005}.

\paragraph{No new theory.} This paper does not prove anything about cohomology, codes, or
contextuality that was not already known. The distance-one statement is an instance of a standard
fact about cycle-space codes; the contribution here is that the audited artifact measures it and
that the paper states it as the repository's own caveat should have.

\paragraph{One machine, one interpreter.} All numbers come from a single Windows 11 machine with
CPython 3.14.6 and NumPy 2.5.1. Every experiment reproduced here is deterministic, so machine
dependence is limited to floating-point details in the Hopfield baseline; the qualitative
comparison does not depend on the hardware, and the third decimal of the printed rates does.

\paragraph{The seeding issue cuts both ways.} The deterministic seeding makes the reproduction
strong and the inference weak: a single stream means the reported rates are one draw from a
distribution whose width the code never reports. This paper attaches a binomial standard error to
one row to make the size of that width visible. It does not re-run the Monte Carlo across seeds,
which would be the right way to publish a rate.

\section{Reproducibility statement}

The package pins the audited commit (\texttt{subjects.lock}), declares every run
(\texttt{experiments/\allowbreak{}plan.toml}), keeps each process's stdout and stderr (\texttt{evidence/\allowbreak{}raw/\allowbreak{}}),
records the byte-identity verdict per run, and derives every number in this text from
\texttt{data/\allowbreak{}measurements.json} through the ledger. The reference list is generated from a curated
spec after a machine resolved each DOI; \texttt{papertool cites} and \texttt{papertool numbers}
fail the build on drift. \texttt{bash verify\_\allowbreak{}all.sh} runs all of it and refuses to pass while any
check fails.

\section{Conclusion}

The audited repository asked whether paradox can be a foundation, answered no in its own README,
and shipped the experiments that support the no. This paper confirms that answer from outside: the
experiments reproduce byte for byte, the contextuality identification is correct, the refutation
of the author's diagonal lemma is correct, and the invariant that survives is a parity checksum
whose minimum distance is one. The one place the audit improves on the text is the third claim,
where "a checksum, not a memory" becomes a measurable statement about distance and undetectable
errors rather than a caveat in a parenthesis.

That is a modest outcome for a large programme, and it is the outcome the programme predicted for
itself. Publishing a prediction like that, and the experiment that fulfils it, is rarer than
publishing the programme.

\section*{Acknowledgements and disclosure}

The author thanks the maintainer of \texttt{paradox-logic} for committing result files next to the
code that produces them, which is what made a byte-level reproduction possible. AI assistance: the
reproduction harness, the plan, the number ledger, the reference verification and the drafting were
produced in a session with a Hermes Agent instance acting as a research assistant. Every number in
this paper is generated by the artifact from raw logs and every reference was resolved by a machine
query before it was cited; the checks that enforce both are in \texttt{verify\_\allowbreak{}all.sh} and were run
by the human author on the published version.

\bibliographystyle{plainnat}
\bibliography{references}

\appendix
\section{Commands}

\begin{verbatim}
bash scripts/fetch_subjects.sh        # clone at the pinned commit
papertool measure --root .            # re-run the plan, keep logs, compare bytes
papertool measure --root . --derive-only
papertool ledger --root .             # measurements -> macros, ledger, figures
papertool bib --root .                # curated sources -> references.bib
bash build.sh                         # all of the above, then the PDF
bash verify_all.sh                    # everything, plus the freshness checks
papertool measure --root .            # run the plan, keep logs
papertool cites --root . --check
papertool numbers --root .
\end{verbatim}

\end{document}
